% GECToR一轮推断：源词元（含句首位置）经编码器得到逐词向量，检测头与标签头逐位置输出，
% 加置信度偏置后取标签并同时应用，结果进入下一轮。概率为教学示例，与正文算例一致。
\begin{tikzpicture}[
  x=1mm,y=1mm,figure picture,
  tok/.style={figure node,minimum width=17mm,minimum height=6mm,
    inner sep=1pt,draw=FigureInput,fill=FigureInputFill,font=\figurefont\scriptsize},
  enc/.style={figure node,minimum height=7mm,inner sep=1.5pt,
    draw=FigureModel,fill=FigureModelFill,font=\figurefont\scriptsize},
  vec/.style={font=\figurefont\scriptsize},
  outc/.style={figure node,minimum width=19.5mm,minimum height=10mm,align=center,
    inner sep=1pt,draw=FigureOutput,fill=FigureOutputFill,font=\figurefont\scriptsize},
  res/.style={figure node,minimum height=6.5mm,inner sep=1.5pt,
    draw=FigureOutput,fill=FigureOutputFill,font=\figurefont\scriptsize},
  note/.style={font=\figurefont\scriptsize,text=FigureMuted,align=center},
  row label/.style={font=\figurefont\scriptsize,text=FigureMuted,anchor=east,align=right},
  flow/.style={figure flow}
]
  \foreach \i/\w/\tag/\pe in {0/{\$START}/KEEP/0.02,1/She/KEEP/0.03,2/go/{VB$\to$VBZ}/0.85,
      3/to/{APPEND\_the}/0.70,4/meeting/KEEP/0.45}{
    \node[tok] (x\i) at (\i*22,0) {\w};
    \node[vec] (h\i) at (\i*22,22) {$\bm h_{\i}$};
    \node[outc] (t\i) at (\i*22,36) {标签：\tag\\有错概率\pe};
    \draw[flow] (x\i) -- (x\i |- 0,8.5);
    \draw[flow] (\i*22,15.5) -- (h\i);
    \draw[flow] (h\i) -- (t\i);
  }
  \node[enc,minimum width=108mm] (enc) at (44,12)
    {预训练编码器（每个词取首个子词的向量）};
  \node[row label] at (-12,0) {本轮输入};
  \node[row label] at (-12,22) {逐词向量};
  \node[row label] at (-12,36) {标签头\\检测头};

  \node[res,minimum width=60mm] (y) at (44,54) {She goes to the meeting};
  \draw[flow] (t2.north) -- (t2.north |- y.south);
  \draw[flow] (t3.north) -- (t3.north |- y.south);
  \node[note,anchor=west] at (76,54) {同时应用，\\作为下一轮输入};
  \node[row label] at (-12,54) {重建结果};
  \node[note] at (88,44.5) {加偏置$\beta$后KEEP胜出};
\end{tikzpicture}
